\subsection{Valuation rings in an algebraic extension}

PROPOSITION 6. Let K be a field, v a valuation on K, A its ring, L an algebraic extension of K and A' the integral closure of A in L. Let V be the set of valuation rings on L which extend v and M' the set of maximal ideals of A'. Then the mapping V \(\mapsto\) m(V) \(\cap\) A' is a bijection of V onto M' and \(m' \mapsto A'_{m'}\) is the inverse bijection.

Every maximal ideal \(m'\) of \(A'\) is such that \(m' \cap A\) is the maximal ideal of A (Chapter V, §2, no. 1, Proposition 1) and \(A_{m'}'\) is dominated by a valuation ring V of L (which is therefore the ring of a valuation on L extending v) (§1, no. 2, Corollary to Theorem 2). The field L is the union of a directed family of sub-extensions \(K_{\alpha}\) of L which are finite over K and it will suffice, in order to see that \(V = A_{m'}'\) , to prove that \(V \cap K_{\alpha} = A_{m'}' \cap K_{\alpha}\) for all \(\alpha\) . Now, if we write \(A_{\alpha}' = A' \cap K_{\alpha}\) , \(A_{\alpha}'\) is the integral closure of A in \(K_{\alpha}\) and hence is the intersection of the rings of the valuations on \(K_{\alpha}\) which extend v and these rings \(V_{i\alpha}\) are finite in number and are the local rings \((\mathrm{A}_{\alpha}')_{\mathrm{m}'_{i\alpha}}\) of \(A_{\alpha}' (1 \leqslant i \leqslant n)\) , where the \(m'_{i\alpha}\) are the distinct maximal ideals of \(A_{\alpha}' (no. 3, Remark)\) ; but \(m' \cap A_{\alpha}'\) is one of the \(m'_{i\alpha}\) and \(V \cap K_{\alpha}\) is therefore equal to the corresponding local ring \((\mathrm{A}'_{\alpha})_{m'_{i\alpha}} \subset \mathrm{A}'_{m'}\) , which completes the proof that \(V = A'_{m'}\) . Conversely, if \(V \in V\) , then \(A' \subset V\) (§3, no. 3, Proposition 6) and, if \(m' = m(V) \cap A'\) , then \(m' \cap A = m\) , hence \(m'\) is a maximal ideal of \(A'\) (Chapter V, §2, no. 1, Proposition 1) and the above argument shows that \(V = A'_{m'}\) .

PROPOSITION 7. Let K be a field, L a quasi-Galois extension of K and f and \(f'\) places of L with values in the same field F. Suppose that the restrictions of f and \(f'\) to K coincide. Then there exists a K-automorphism s of L such that \(f' = f \circ s\) .

Let A be the ring of the place of K the common restriction of f and \(f'\) . The rings of f and \(f'\) contain the integral closure \(A'\) of A in L (§ 1, no. 3, Corollary 3 to Theorem 3) and hence (Chapter V, § 2, no. 3, Corollary 1 to Proposition 6) there exists a K-automorphism s of L such that the restrictions of \(f'\) and \(f \circ s\) to \(A'\) are equal; if \(m'\) is the common kernel of these restrictions, \(m' \cap A\) is the maximal ideal of A, hence \(m'\) is a maximal ideal of \(A'\) and the places \(f'\) and \(f \circ s\) coincide on the ring \(A_{m'}\) ; but by Proposition 6 the only valuation ring of L dominating \(A_{m'}\) is the ring \(A_{m'}\) itself and hence the rings of the places \(f'\) and \(f \circ s\) are the same.

COROLLARY 1. Let K be a field, v a valuation on K, L a quasi-Galois extension of K and v' and v'' two extensions of v to L. Then there exists a K-automorphism s of L such that v'' is equivalent to v' o s.

Let \(f'\) and \(f''\) be the places of K associated with \(v'\) and \(v''\) ; replacing them if need be by equivalent places, it may be assumed that they both take their values in the algebraic closure of the residue field of \(v\) (no. 1, Proposition 1). Then there exists a K-automorphism \(s\) of L such that \(f'' = f' \circ s\) (Proposition 7); thus \(v''\) is equivalent to \(v' \circ s\) by virtue of the correspondence between places and valuations (§ 3, no. 3).

COROLLARY 2. Let K be a field, f a place of K (resp. v a valuation on K) and L a radicial extension of K. Then all the extensions of f (resp. v) to L are equivalent.

L is a quasi-Galois extension and its only automorphism is the identity. Corollary 2 therefore follows from Proposition 7 (resp. Corollary 1).

PROPOSITION 8. Let K be a field, v a valuation on K, L a finite quasi-Galois extension of K of degree n and \((v_{i}^{\prime})_{1\leq i\leq g}\) a complete system of extensions of v to L. Then \(e(v_{i}^{\prime}|v)\) and \(f(v_{i}^{\prime}|v)\) have values e and f independent of i. Then \(efg \leqslant n\) . If the integral closure in L of the ring A of v is a finitely generated A-module, then \(efg = n\) .

This follows immediately from Theorems 1 (no. 3) and 2 (no. 5).

